\par\medskip
\noindent
\underline{Conjecture \SGARefTarget{sga7.sga7i.eix.n14-4}{14.4}.}  The hypotheses being those just set out in
\SGARefLink{sga7.sga7i.eix.n14-3-5}{14.3.5}, there exists a homomorphism

\begin{equation}\tag{\SGARefTarget{sga7.sga7i.eix.n14-4-1}{14.4.1}}
 i:\underline M_U\longrightarrow G_U
\end{equation}

\noindent
such that, for every prime number $\ell$, the extension (\SGARefLink{sga7.sga7i.eix.n14-3-5-6}{14.3.5.6}) is
isomorphic to the extension deduced from $i$ as explained for (\SGARefLink{sga7.sga7i.eix.n14-2-2}{14.2.2}).
More precisely, there is one and only one way of associating with every
situation of the type considered in \SGARefLink{sga7.sga7i.eix.n14-3-5}{14.3.5} a homomorphism (\SGARefLink{sga7.sga7i.eix.n14-4-1}{14.4.1})
satisfying the preceding condition, in such a way that the formation of
$i$ is compatible with every base change $S'\longrightarrow S$
associated with a continuous homomorphism $V\longrightarrow V'$ of
noetherian adic rings.

\medskip
\noindent
\SGARefTarget{sga7.sga7i.eix.n14-4-2}{14.4.2}. It seems that the existence of $i$, compatibly with base change,
follows from M. Raynaud's rigid-analytic theory, in view of the fact that
he is developing it in a framework more general than that stated in
\SGARefLink{sga7.sga7i.eix.n14-1}{14.1}, working, as here, over general noetherian adic rings.  Moreover,
to prove the existence and uniqueness asserted in \SGARefLink{sga7.sga7i.eix.n14-4}{14.4}, one is
immediately reduced to considering only those $S$ for which $S_o$ is a
scheme of finite type over $\mathbb Z$.  Indeed, $V$ is the filtered
inductive limit of subrings $V_i$ of finite type over $\mathbb Z$; these
may be supposed normal by Nagata's theorem (EGA IV 7.8), and, for $i$
sufficiently large, the situation $(A,A',W)$ comes from an analogous
situation over $V_i$, or alternatively over the separated completion of
$V_i$ for the topology defined by $m_i=\underline m V_i$.  This gives
the asserted reduction.

\medskip
\noindent
\SGARefTarget{sga7.sga7i.eix.n14-4-3}{14.4.3}. Observe now that $S-S_o$ is a Jacobson scheme (EGA IV 10.5.8).
Thus, to see that two solutions $i,i'$ of the problem, for a given
situation $(A,A',W)$, are equal, it is enough to prove that these
homomorphisms are equal at every \underline{closed} point $\eta$ of
$U$.  Let $S'$ be the normalization of the closure of $\eta$ in $S$.
Then $S'$ is a complete trait, and the uniqueness result of \SGARefLink{sga7.sga7i.eix.n4-3}{4.3} may be
applied to the pullback to $S'$ of the situation over $S$, since the
residue field of $S'$ is \underline{finite}.  Once uniqueness is thus
established, it implies compatibility of the formation of $i$ with base
change, subject to existence.

\medskip
\noindent
\underline{Remarks \SGARefTarget{sga7.sga7i.eix.n14-5}{14.5}.}  It would remain to formulate with the desired
generality, and to establish, a compatibility between the homomorphism
$i$ introduced in \SGARefLink{sga7.sga7i.eix.n14-4}{14.4} by a rigid-analytic method and the analogous
homomorphism (\SGARefLink{sga7.sga7i.eix.n13-5-13}{13.5.13}) of complex-analytic theory---or, better still, to
find a common generalization of them!

